\documentclass{article}
\usepackage{graphicx} % Required for inserting images
\usepackage{titling}
\newcommand{\subtitle}[1]{%
  \posttitle{%
    \par\end{center}
    \begin{center}\large#1\end{center}
    \vskip0.5em}%
}
\usepackage{geometry}
 \geometry{
 a4paper,
 left=30mm,
 right=20mm,
 top=25mm,
 bottom=25mm,
 }

\title{Title}
\subtitle{(Professor)}
\author{Author}
\date{Date}

\begin{document}

\maketitle

\section*{Presentation summary}
En esta sección va un resumen de la charla presentada por el profesor. Para citar un artículo tienen que agregar el texto extraído del bib en el archivo \textbf{mybib.bib} y citarlo como se muestra a continuación \cite{Tahamtan2016}. La citación es automática, si quieren hacer múltiples citas pueden agregar las referencias separadas con una coma \cite{Tahamtan2016, Kurtz2005}. Aunque no es necesario podrían agregar ecuaciones, para esto pueden usar el comando \textbf{equation} como se muestra a continuación en (\ref{Ecuacion1}). Notar que las ecuaciones se pueden referenciar dentro del texto utilizando el comando \textbf{ref} y agregando la etiqueta en la ecuación mediante el comando \textbf{label}.

\begin{equation}\label{Ecuacion1}
    E = mc^2.
\end{equation}

De la misma forma pueden citar las tablas utilizando el mismo comando anterior, la Tabla \ref{Mitabla} es un ejemplo de cómo hacer una tabla random.

\begin{table}[h]
    \centering
    \begin{tabular}{cc}
    \hline
        Esto es & una tabla \\
    \hline
        Continuación & \dots \\
        Continuación & \dots \\
        \hline
    \end{tabular}
    \caption{Texto que va bajo la tabla}
    \label{Mitabla}
\end{table}

Finalmente podemos citar las secciones, en la sección \ref{Fortalezas} veremos las fortalezas de la investigación. En la sección \ref{Debilidades} mencionamos algunas debilidades encontradas de la investigación. La sección de preguntas no está enumerada, esto lo podemos hacer agregando un {*} para que en el texto no se agregue un número.

\section{Strengths}\label{Fortalezas}

Acá tienen que agregar las fortalezas de la investigación.

\section{Weakness}\label{Debilidades}

Acá agregan debilidades o comentarios respecto a la investigación.

\section*{Questions}

Acá preguntas que se podrían abordar en otras investigaciones o cuestiones respecto a la metodología o resultados.

\bibliographystyle{plain}
\bibliography{mybib.bib}

\end{document}
